\chapter{Lesson 4: Greetings and Classroom Expressions}

\section{Everyday greetings}

Greetings (\ja{あいさつ}, \emph{aisatsu}) are fixed expressions. Many have no literal
equivalent in English, so learn them as whole units tied to a situation.

\begin{center}
\begin{tabular}{lll}
\toprule
Japanese & Romaji & When / meaning \\
\midrule
\ja{おはようございます} & ohayō gozaimasu & good morning (polite) \\
\ja{おはよう} & ohayō & good morning (to friends, family) \\
\ja{こんにちは} & konnichiwa & hello, good afternoon \\
\ja{こんばんは} & konbanwa & good evening \\
\ja{さようなら} & sayōnara & goodbye (for a long time; teacher at end of class) \\
\ja{じゃあ、また} & jā, mata & see you (casual) \\
\ja{おやすみなさい} & oyasumi nasai & good night (before sleep) \\
\ja{ありがとうございます} & arigatō gozaimasu & thank you (polite) \\
\ja{どうも} & dōmo & thanks (short; also ``hi'') \\
\ja{いいえ} & iie & no; ``not at all'' (reply to thanks) \\
\ja{すみません} & sumimasen & excuse me; I'm sorry; (sometimes) thank you \\
\ja{ごめんなさい} & gomen nasai & I'm sorry (personal apology) \\
\ja{はじめまして} & hajimemashite & nice to meet you (first meeting only) \\
\ja{よろしくおねがいします} & yoroshiku onegai shimasu & please treat me well \\
\bottomrule
\end{tabular}
\end{center}

\begin{remark}
The \ja{は} at the end of \ja{こんにちは} and \ja{こんばんは} is the topic particle and
is read \emph{wa}. These greetings are fragments of old sentences: ``As for today
(\ja{こんにち}) \ldots''. Writing \ja{こんにちわ} is a common mistake (and is marked
wrong on quizzes).
\end{remark}

\subsection{Mealtime and home expressions}
\begin{center}
\begin{tabular}{lll}
\toprule
Japanese & Romaji & Use \\
\midrule
\ja{いただきます} & itadakimasu & before eating (``I humbly receive'') \\
\ja{ごちそうさまでした} & gochisōsama deshita & after eating (thanks for the meal) \\
\ja{いってきます} & ittekimasu & leaving home (``I'll go and come back'') \\
\ja{いってらっしゃい} & itterasshai & reply: ``go and come back'' \\
\ja{ただいま} & tadaima & arriving home: ``I'm home'' \\
\ja{おかえりなさい} & okaerinasai & reply: ``welcome back'' \\
\bottomrule
\end{tabular}
\end{center}

\begin{center}
\begin{tikzpicture}[sp/.style={draw, rounded corners, fill=blue!6, minimum width=3.4cm, minimum height=0.8cm},
    re/.style={draw, rounded corners, fill=orange!10, minimum width=3.4cm, minimum height=0.8cm}]
  \node[sp] (a1) at (0,0) {\ja{いってきます}};
  \node[re] (b1) at (5.5,0) {\ja{いってらっしゃい}};
  \node[sp] (a2) at (0,-1.4) {\ja{ただいま}};
  \node[re] (b2) at (5.5,-1.4) {\ja{おかえりなさい}};
  \node[sp] (a3) at (0,-2.8) {\ja{ありがとうございます}};
  \node[re] (b3) at (5.5,-2.8) {\ja{いいえ}};
  \node[sp] (a4) at (0,-4.2) {\ja{はじめまして}};
  \node[re] (b4) at (5.5,-4.2) {\ja{はじめまして}};
  \foreach \i in {1,2,3,4} \draw[-Stealth, thick] (a\i) -- (b\i);
  \node[font=\small\bfseries] at (0,0.8) {first speaker};
  \node[font=\small\bfseries] at (5.5,0.8) {standard reply};
\end{tikzpicture}
\end{center}

\section{Classroom expressions}

\subsection{What the teacher says}
\begin{center}
\begin{tabular}{lll}
\toprule
Japanese & Romaji & Meaning \\
\midrule
\ja{きいてください} & kiite kudasai & please listen \\
\ja{みてください} & mite kudasai & please look \\
\ja{いってください} & itte kudasai & please say it \\
\ja{かいてください} & kaite kudasai & please write \\
\ja{よんでください} & yonde kudasai & please read \\
\ja{もういちど} & mō ichido & one more time \\
\ja{ゆっくり} & yukkuri & slowly \\
\ja{しゅくだい} & shukudai & homework \\
\ja{しけん} & shiken & exam \\
\ja{しつもん} & shitsumon & question \\
\bottomrule
\end{tabular}
\end{center}
(We will learn in Lesson 11 that \ja{〜てください} is the polite request form.)

\subsection{What students say}
\begin{center}
\begin{tabular}{lll}
\toprule
Japanese & Romaji & Meaning \\
\midrule
\ja{わかりました} & wakarimashita & I understood \\
\ja{わかりません} & wakarimasen & I don't understand / know \\
\ja{しりません} & shirimasen & I don't know (a fact) \\
\ja{ちょっとまってください} & chotto matte kudasai & please wait a moment \\
\ja{もういちどおねがいします} & mō ichido onegai shimasu & once more, please \\
\ja{ゆっくりおねがいします} & yukkuri onegai shimasu & more slowly, please \\
\ja{〜はにほんごでなんですか} & \ldots wa Nihongo de nan desu ka & how do you say \ldots\ in Japanese? \\
\ja{〜ってどういういみですか} & \ldots tte dō iu imi desu ka & what does \ldots\ mean? \\
\bottomrule
\end{tabular}
\end{center}

\begin{example}[Asking for a word]
\leavevmode
\begin{itemize}
  \item[A:] \ja{すみません。「pencil」はにほんごでなんですか。}\\
        \emph{Sumimasen. ``Pencil'' wa Nihongo de nan desu ka.}\\
        Excuse me. How do you say ``pencil'' in Japanese?
  \item[B:] \ja{「えんぴつ」です。} \emph{``Enpitsu'' desu.} It's ``enpitsu''.
  \item[A:] \ja{ありがとうございます。} \emph{Arigatō gozaimasu.} Thank you.
\end{itemize}
\end{example}

\section{Politeness: a first look}

Japanese speech changes depending on who you are talking to. Two basic levels appear
right away:
\begin{itemize}
  \item \textbf{polite} (\ja{です}/\ja{ます} style): to teachers, strangers, older people,
        in class. This is the default for this course.
  \item \textbf{casual/plain}: to close friends, family, younger people. Shorter greetings
        like \ja{おはよう}, \ja{ありがとう}, \ja{じゃあね}.
\end{itemize}
Adding \ja{ございます} raises \ja{おはよう} and \ja{ありがとう} to polite. Using
casual forms with a teacher sounds rude; using very polite forms with a close friend
sounds cold.

\begin{aside}
\textbf{Keigo (\ja{けいご}) awareness.} Beyond plain vs.\ polite, Japanese has an
elaborate system of honorific language: \emph{sonkeigo} (raising the other person),
\emph{kenjōgo} (humbling oneself) and \emph{teineigo} (the polite \ja{です}/\ja{ます}
style we are learning). Shop clerks greet customers with \ja{いらっしゃいませ}
(``welcome'') — a very respectful form one does not reply to; a nod is enough. We will
not produce sonkeigo or kenjōgo this term, but we will hear them in shops and
announcements.
\end{aside}

\section{Bowing}

Bowing (\ja{おじぎ}, \emph{ojigi}) accompanies greetings, thanks and apologies. The depth
of the bow reflects the situation.

\begin{center}
\begin{tikzpicture}[x=1cm,y=1cm]
  \foreach \ang/\name/\use [count=\i] in {15/eshaku/casual greeting, 30/keirei/to customers or superiors, 45/saikeirei/deep thanks or apology} {
    \begin{scope}[xshift=\i*4.2cm-4.2cm]
      \draw[gray] (-0.6,0) -- (1.8,0);
      \draw[very thick] (0,0) -- (0,1.5);           % legs
      \draw[very thick] (0,1.5) -- ++(90-\ang:1.5) coordinate (n); % torso
      \draw[thick, fill=white] ($(n)+(90-\ang:0.3)$) circle (0.3);   % head
      \draw[dashed, gray] (0,1.5) -- (0,3.4);
      \draw[-Stealth] (0,2.6) arc (90:90-\ang:1.1);
      \node[font=\small] at (0.6,-0.4) {\name\ ($\ang^\circ$)};
      \node[font=\scriptsize, text width=3.4cm, align=center] at (0.6,-0.9) {\use};
    \end{scope}
  }
\end{tikzpicture}
\end{center}

\begin{itemize}
  \item Men usually keep their hands at their sides; women often place hands together
        in front.
  \item Keep your eyes down; do not bow while staring at the other person.
  \item Between people of different status, the junior bows lower and longer.
  \item A bow can replace a handshake, though many Japanese businesspeople will offer
        a handshake to foreigners. A small head nod (\emph{eshaku}) is fine with friends.
\end{itemize}

\section{Names and titles}

\begin{important}[Name order and \ja{さん}]
\begin{itemize}
  \item Japanese names are given \textbf{family name first}: \ja{やまだ　はなこ}
        \emph{Yamada Hanako} (Yamada = family name).
  \item Attach \ja{さん} \emph{-san} to \emph{other people's} names (family or given):
        \ja{やまださん}. \textbf{Never} use \ja{さん} with your own name.
  \item Teachers are addressed as \ja{せんせい} \emph{sensei}, either alone or after the
        family name: \ja{やまだせんせい}. Do not call a teacher ``\ja{あなた}'' (you).
  \item Other suffixes: \ja{くん} \emph{-kun} (boys, juniors), \ja{ちゃん} \emph{-chan}
        (children, close friends), \ja{さま} \emph{-sama} (very formal; customers).
\end{itemize}
\end{important}

\begin{aside}
In Japan people use family names far more than in Canada; classmates may call each
other by family name + \ja{さん} for the whole year. Calling someone by their given name
with no suffix (\emph{yobisute}) is reserved for close friends and family. When in
doubt, family name + \ja{さん}.
\end{aside}

\section*{Lesson 4 review}
\begin{itemize}
  \item Greetings depend on time of day and relationship: \ja{おはようございます},
        \ja{こんにちは}, \ja{こんばんは}, \ja{さようなら}, \ja{おやすみなさい}.
  \item \ja{すみません} = excuse me / sorry / (sometimes) thanks.
  \item Polite (\ja{です}/\ja{ます}) style is our default; keigo exists at higher levels.
  \item Bow depth reflects formality. Family name first; \ja{さん} for others only.
\end{itemize}
